
\documentclass{article}
\usepackage{colortbl}
\usepackage{makecell}
\usepackage{multirow}
\usepackage{supertabular}

\begin{document}

\newcounter{utterance}

\centering \large Interaction Transcript for game `dond', experiment `semi\_en', episode 12 with reflexion{-}Qwen3.5{-}27B.
\vspace{24pt}

{ \footnotesize  \setcounter{utterance}{1}
\setlength{\tabcolsep}{0pt}
\begin{supertabular}{c@{$\;$}|p{.15\linewidth}@{}p{.15\linewidth}p{.15\linewidth}p{.15\linewidth}p{.15\linewidth}p{.15\linewidth}}
    \# & $\;$A & \multicolumn{4}{c}{Game Master} & $\;\:$B\\
    \hline

    \theutterance \stepcounter{utterance}  
    & & \multicolumn{4}{p{0.6\linewidth}}{
        \cellcolor[rgb]{0.9,0.9,0.9}{
            \makecell[{{p{\linewidth}}}]{
                \texttt{\tiny{[P1$\langle$GM]}}
                \texttt{You have attempted this same task before and did not succeed. These are the plans you wrote after those attempts.} \\
\texttt{Trial \#1: New plan: I failed because I submitted a proposal immediately after the other player did, without engaging in the negotiation phase to understand their value function or confirm a complementary split. By stating my preference for the nut and offering the rest, I should have waited for their response to see if they actually wanted the shells and clocks. If they had proposed taking the shells and clocks, I could have confirmed the deal. Instead, I assumed the split was agreed upon and submitted a proposal that likely conflicted with theirs (e.g., if they also wanted the nut or didn't want the shells/clocks). Next time, I will use my first message to propose a split, wait for the other player's counter{-}proposal or confirmation, and only submit a secret proposal once I am certain our desired items do not overlap. Specifically, I will ask, "Do you want the 2 shells and 2 clocks in exchange for the 1 nut?" and wait for their reply before finalizing.} \\
\\ 
\texttt{Trial \#2: New plan: I failed because I submitted a secret proposal immediately after the other player verbally agreed, without verifying that their *secret* proposal would match my verbal agreement. The other player stated they value clocks highly, yet I offered them the clocks in my verbal deal while they likely intended to claim the nut (which I also wanted) or the clocks in their secret proposal, causing a conflict. Next time, after the other player verbally agrees to a split, I will explicitly ask them to confirm their intended secret proposal format (e.g., "To ensure we don't clash, will your secret proposal be '{[}Proposal: 2 shells, 2 clocks{]}'?"). Only after receiving a clear confirmation that their desired items are exactly the ones I am leaving for them will I submit my own secret proposal. If they do not explicitly confirm the specific items they will claim, I will not submit a proposal yet to avoid a zero{-}score outcome.} \\
\\ 
\\ 
\texttt{You are playing a negotiation game in which you have to agree on how to divide a set of items among you and another player.} \\
\\ 
\texttt{Rules:} \\
\texttt{(a) You and the other player are given a set of items. Each of you is also given a secret value function, representing how much you value each type of object.} \\
\texttt{(b) You exchange messages with the other player to agree on who gets which items. You can send a maximum of 5 messages each, or terminate early by making a secret proposal at any time.} \\
\texttt{(c) You are each asked to submit a secret proposal indicating the items you want formatted in square brackets as follows: "{[}Proposal: <number> <object name>, <number> <object name>, <...>{]}"} \\
\texttt{(d) If your proposals are complementary, i.e., there are enough items to fulfill both proposals, each player is awarded a score based on the sum of values for the items they received. Otherwise, both of you get zero points.} \\
\texttt{(e) Your goal is to maximize the score you receive.} \\
\\ 
\texttt{Let us start.} \\
\\ 
\texttt{The set of available items is:} \\
\\ 
\texttt{2 shells, 1 nut, 2 clocks} \\
\\ 
\texttt{Your secret value function is:} \\
\\ 
\texttt{shell: 1, nut: 8, clock: 0} \\
\\ 
\texttt{IMPORTANT: Your messages, unless it is the secret proposal, are directly transmitted to the other player, so do not include any response to the rules or text announcing your message. To make a secret proposal, use the indicated format. Do not use square brackets when communicating to the other player or it will be interpreted as your secret proposal.} \\
            }
        }
    }
    & \\ \\

    \theutterance \stepcounter{utterance}  
    & \multicolumn{4}{p{0.6\linewidth}}{
        \cellcolor[rgb]{0.9,0.9,0.9}{
            \makecell[{{p{\linewidth}}}]{
                \texttt{\tiny{[P1$\rangle$GM]}}
                \texttt{I value the nut highly (8 points) and have no use for the clocks (0 points) or shells (1 point each). To maximize our combined score and ensure we don't get zero points, I propose a split where I take the **1 nut**, and you take the **2 shells and 2 clocks**.} \\
\\ 
\texttt{Does this work for you? If you agree, please confirm that your secret proposal will be exactly: `{[}Proposal: 2 shells, 2 clocks{]}`. Once you confirm, I will submit my proposal for the nut.} \\
            }
        }
    }
    & & \\ \\

    \theutterance \stepcounter{utterance}  
    & & & \multicolumn{2}{p{0.3\linewidth}}{
        \cellcolor[rgb]{0.9,0.9,0.9}{
            \makecell[{{p{\linewidth}}}]{
                \texttt{\tiny{[GM$|$GM]}}
                \texttt{proposal} \\
            }
        }
    }
    & & \\ \\

    \theutterance \stepcounter{utterance}  
    & & & \multicolumn{2}{p{0.3\linewidth}}{
        \cellcolor[rgb]{0.9,0.9,0.9}{
            \makecell[{{p{\linewidth}}}]{
                \texttt{\tiny{[GM$|$GM]}}
                [2, 0, 2]
            }
        }
    }
    & & \\ \\

    \theutterance \stepcounter{utterance}  
    & & \multicolumn{4}{p{0.6\linewidth}}{
        \cellcolor[rgb]{0.9,0.9,0.9}{
            \makecell[{{p{\linewidth}}}]{
                \texttt{\tiny{[GM$\rangle$P2]}}
                \texttt{You have attempted this same task before and did not succeed. These are the plans you wrote after those attempts.} \\
\texttt{Trial \#1: New plan: I failed because I explicitly revealed my secret value function (stating I value clocks at 3) in my public message, which allowed the other player to exploit my preferences, and I likely misinterpreted the negotiation outcome or the other player's hidden values. My proposal of taking only the nut was also inconsistent with my stated preference for clocks in the previous turn, causing confusion or a non{-}complementary agreement. Next time, I will strictly keep my value function secret in all public messages, only stating my desired items without explaining why. I will negotiate to secure the nut (my highest value) and the clocks (my second highest) by offering the shells (worth 0 to me) to the other player. I will ensure my final secret proposal matches the agreed{-}upon split exactly, likely proposing "{[}Proposal: 1 nut, 2 clocks{]}" if the other player agrees to take the shells, or adjusting based on their counter{-}offer while maintaining the goal of maximizing my score.} \\
\\ 
\texttt{Trial \#2: New plan: I failed because I misread the negotiation outcome and submitted a proposal that did not align with the agreement I just verbally confirmed. In my public message, I explicitly agreed to take the "2 shells and 2 clocks" while the other player took the "1 nut," but my secret proposal requested only the "1 nut." This created a conflict where both players claimed the nut, resulting in zero points. Next time, I will rigorously cross{-}reference my final secret proposal with the specific items I verbally agreed to receive in the last message. If I agree to take shells and clocks, my proposal must strictly list "{[}Proposal: 2 shells, 2 clocks{]}" and exclude the nut, ensuring the proposals are complementary. I will also avoid revealing my secret value function in public messages to prevent the other player from exploiting my preferences.} \\
\\ 
\\ 
\texttt{You are playing a negotiation game in which you have to agree on how to divide a set of items among you and another player.} \\
\\ 
\texttt{Rules:} \\
\texttt{(a) You and the other player are given a set of items. Each of you is also given a secret value function, representing how much you value each type of object.} \\
\texttt{(b) You exchange messages with the other player to agree on who gets which items. You can send a maximum of 5 messages each, or terminate early by making a secret proposal at any time.} \\
\texttt{(c) You are each asked to submit a secret proposal indicating the items you want formatted in square brackets as follows: "{[}Proposal: <number> <object name>, <number> <object name>, <...>{]}"} \\
\texttt{(d) If your proposals are complementary, i.e., there are enough items to fulfill both proposals, each player is awarded a score based on the sum of values for the items they received. Otherwise, both of you get zero points.} \\
\texttt{(e) Your goal is to maximize the score you receive.} \\
\\ 
\texttt{Let us start.} \\
\\ 
\texttt{The set of available items is:} \\
\\ 
\texttt{2 shells, 1 nut, 2 clocks} \\
\\ 
\texttt{Your secret value function is:} \\
\\ 
\texttt{shell: 0, nut: 4, clock: 3} \\
\\ 
\texttt{IMPORTANT: Your messages, unless it is the secret proposal, are directly transmitted to the other player, so do not include any response to the rules or text announcing your message. To make a secret proposal, use the indicated format. Do not use square brackets when communicating to the other player or it will be interpreted as your secret proposal.} \\
\\ 
\\ 
\texttt{The other player has submitted a secret proposal. It is now your turn to submit a proposal of your own.} \\
            }
        }
    }
    & \\ \\

    \theutterance \stepcounter{utterance}  
    & & & \multicolumn{4}{p{0.6\linewidth}}{
        \cellcolor[rgb]{0.9,0.9,0.9}{
            \makecell[{{p{\linewidth}}}]{
                \texttt{\tiny{[GM$\langle$P2]}}
                \texttt{{[}Proposal: 1 nut, 2 clocks{]}} \\
            }
        }
    }
     \\ \\

    \theutterance \stepcounter{utterance}  
    & & & \multicolumn{2}{p{0.3\linewidth}}{
        \cellcolor[rgb]{0.9,0.9,0.9}{
            \makecell[{{p{\linewidth}}}]{
                \texttt{\tiny{[GM$|$GM]}}
                \texttt{proposal} \\
            }
        }
    }
    & & \\ \\

    \theutterance \stepcounter{utterance}  
    & & & \multicolumn{2}{p{0.3\linewidth}}{
        \cellcolor[rgb]{0.9,0.9,0.9}{
            \makecell[{{p{\linewidth}}}]{
                \texttt{\tiny{[GM$|$GM]}}
                [0, 1, 2]
            }
        }
    }
    & & \\ \\

    \theutterance \stepcounter{utterance}  
    & & & \multicolumn{2}{p{0.3\linewidth}}{
        \cellcolor[rgb]{0.9,0.9,0.9}{
            \makecell[{{p{\linewidth}}}]{
                \texttt{\tiny{[GM$|$GM]}}
                [[2, 0, 2], [0, 1, 2]]
            }
        }
    }
    & & \\ \\

    \theutterance \stepcounter{utterance}  
    & \multicolumn{4}{p{0.6\linewidth}}{
        \cellcolor[rgb]{0.95,0.95,0.95}{
            \makecell[{{p{\linewidth}}}]{
                \texttt{\tiny{[P2$\rightarrow$P1]}}
                \texttt{Post{-}episode Reflexion (not sent to GM):} \\
\texttt{New plan: I failed because I attempted to force the other player to confirm a specific secret proposal format within the public chat, which is a strategic error. The rules state that secret proposals are submitted independently and are not visible to the other player; therefore, I cannot verify their actual submission or force them to "confirm" a specific string in the chat. My previous strategy of asking "Does this work?" and waiting for a confirmation of the *format* created a false sense of security. Instead, I will focus on negotiating a mutually beneficial split based on inferred values. Since I value the nut (8) and have no use for clocks (0), I will explicitly offer the clocks and shells to the other player in exchange for the nut, asking simply if they are willing to take those items. I will only submit my secret proposal `{[}Proposal: 1 nut{]}` after the other player explicitly states in the chat that they *want* the shells and clocks (or at least the clocks, since I value them at 0). I will not ask them to confirm the bracketed format, as that is an internal action they cannot verify for me. If they express interest in the nut, I will negotiate a different split or hold off on submitting a proposal until a clear, non{-}overlapping agreement is verbally established.} \\
            }
        }
    }
    & & \\ \\

    \theutterance \stepcounter{utterance}  
    & & & \multicolumn{4}{p{0.6\linewidth}}{
        \cellcolor[rgb]{0.95,0.95,0.95}{
            \makecell[{{p{\linewidth}}}]{
                \texttt{\tiny{[P1$\rightarrow$P2]}}
                \texttt{Post{-}episode Reflexion (not sent to GM):} \\
\texttt{New plan: I failed because I submitted a proposal claiming the nut and both clocks without verifying if the other player's secret proposal left those items available. Since the outcome was a loss, the other player likely claimed the nut or at least one clock in their proposal, creating a conflict. My strategy of greedily taking all high{-}value items (nut: 4, clocks: 3) without a prior public negotiation to establish a complementary split was flawed. Next time, I will first send a public message to negotiate a specific division, explicitly offering the shells (which I value at 0) to the other player in exchange for the nut and clocks. I will only submit a secret proposal after the other player has verbally agreed to take the shells. If I cannot negotiate, I will assume the other player might want the nut and adjust my proposal to claim only the clocks, or vice versa, to ensure the proposals do not overlap. I will strictly ensure my final proposal lists only items the other player has agreed to leave for me.} \\
            }
        }
    }
     \\ \\

\end{supertabular}
}

\end{document}
